\chapter{Retained-Teacher Neural OT Warm Starts}
\label{ch:bf-dpf-learned-ot}

The previous chapter fixed entropic OT with finite Sinkhorn iteration and
barycentric projection as a fully declared teacher object.  The present chapter
studies the narrowest neural extension of that route: do not replace the teacher;
learn how to start it better.

This is the main implementation-bearing neural chapter of the OT block.  Its
question is narrower than ``can a network learn transport?''  It asks instead:
can a network predict a useful internal solver state so that a retained
corrective Sinkhorn phase needs less work while still protecting the declared
finite teacher object?

\section{What problem this chapter solves}
\label{sec:bf-learned-ot-objects}

Let \(X=(x_1,\ldots,x_N)^\top\in\mathbb R^{N\times d_x}\) be the weighted
source cloud and let \(w\in\Delta_N\) be its normalized weights.  Let
\(u=(1/N,\ldots,1/N)\) be the equal target marginal, and let
\(C(X,Z)\in\mathbb R^{N\times N}\) be the declared transport cost matrix.
Chapter~\ref{ch:bf-dpf-entropic-ot-sinkhorn} already told us how to turn this
weighted cloud into an equally weighted cloud by computing a finite
entropic-OT/Sinkhorn teacher and then barycentrically projecting it.

The practical problem now is not to change that object, but to reduce the work
needed to compute it.  So the chapter asks one focused implementation question:
\begin{quote}
Can we learn a solver-relevant latent state that makes the retained finite
Sinkhorn correction cheaper while still returning to the same declared teacher
object after correction?
\end{quote}

\section{Running warm-start cell}
\label{sec:bf-learned-ot-running-cell}

Keep the same three-particle setting as the previous OT chapters: a weighted
source cloud, an equal target marginal, a declared quadratic cost, a
regularization level \(\varepsilon\), and a finite Sinkhorn budget \(K\).  The
finite numerical teacher is the coupling
\(\Pi_{\varepsilon,K}(X,w)\) produced by that declared Sinkhorn route.

A warm-start learner does not try to predict the transported output cloud first.
It tries to predict a latent teacher coordinate, such as dual variables,
log-scalings, or scaling vectors, that initializes the retained corrective solve.
Write
\begin{equation}
\label{eq:bf-learned-ot-student-map}
  h_\psi = S_\psi(X,w;\varepsilon,N,d_x),
\end{equation}
where \(h_\psi\) is later translated into the initial state of the corrective
Sinkhorn loop.

The conceptual contrast is then simple:
\begin{itemize}
  \item a naive initialization changes only how fast the corrective solve reaches
  the teacher,
  \item a learned initialization also changes only that execution path,
  \item the corrected coupling and barycentric output are still supposed to be
  judged against the same declared finite teacher object.
\end{itemize}

\paragraph{Plain-language takeaway.}
A warm-start learner is not asked to invent the transport result.  It is asked to
guess a better place from which the retained teacher solver should begin.

\section{Object audit: what changes and what does not}
\label{sec:bf-learned-ot-hierarchy}

There is no single teacher unless the variant is named.  In this chapter, the
main teacher is the retained EOT/Sinkhorn route from the previous chapter.  The
most important teacher-side objects are:
\begin{enumerate}[label=(\roman*)]
  \item the exact regularized coupling \(\Pi_\varepsilon^\star\),
  \item the finite numerical teacher \(\Pi_{\varepsilon,K}\),
  \item the barycentric teacher output
  \begin{equation}
  \label{eq:bf-learned-ot-barycentric-teacher}
    Y_\tau(X,w,\varepsilon)=B(\Pi_\tau(X,w,\varepsilon)).
  \end{equation}
\end{enumerate}

The student does not directly define a new target cloud at first.  It predicts a
solver-relevant latent state, and the retained corrective Sinkhorn phase then
returns a corrected coupling
\(\widehat\Pi_{\psi,\varepsilon,K_{\mathrm{corr}}}\).

\begin{center}
\small
\begin{tabular}{p{0.25\textwidth} p{0.18\textwidth} p{0.18\textwidth} p{0.27\textwidth}}
\toprule
Quantity & Produced by & Allowed to change? & Deployment role \\
\midrule
latent initialization \(h_\psi\) & student & yes & accelerates execution path \\
finite corrected coupling \(\widehat\Pi_{\psi,\varepsilon,K_{\mathrm{corr}}}\) & retained solve & only through declared correction budget & should match the declared finite teacher closely \\
output cloud \(B(\widehat\Pi_{\psi,\varepsilon,K_{\mathrm{corr}}})\) & barycentric projection & no new object intended & enters the downstream equal-weight layer \\
\bottomrule
\end{tabular}
\end{center}

In the vocabulary of this monograph, this is a retained-teacher or
execution-path approximation.  That is different from routes that change the
transport object itself by learning a direct endpoint map, factorizing the
coupling class, or moving to a richer learned path or operator.  Those routes
may still be scientifically useful, but they answer a different object-audit
question.

\paragraph{Reader checkpoint.}
The retained-teacher claim here is narrow: change the internal execution path,
not the declared corrected transport object.

\section{Algorithm: retained-teacher warm-started Sinkhorn}
\label{sec:bf-learned-ot-teacher}

\begin{algorithm}[ht]
\caption{\texttt{train\_and\_deploy\_warmstarted\_sinkhorn}}
\begin{algorithmic}[1]
\Require training distribution \(\mathcal D\) over weighted particle clouds, declared cost construction, regularization schedule \(\varepsilon\), teacher Sinkhorn solver, architecture class \(S_\psi\), deployment-time corrective budget \(K_{\mathrm{corr}}\)
\Ensure trained predictor \(S_\psi\), corrected deployment coupling \(\widehat\Pi_{\psi,\varepsilon,K_{\mathrm{corr}}}\), barycentric output cloud \(\widetilde X_{\psi,\mathrm{corr}}\), and declared residual diagnostics
\State Generate teacher problems by solving the retained EOT/Sinkhorn route on sampled clouds from \(\mathcal D\)
\State Extract the latent teacher target to be predicted, e.g. dual/scaling coordinates
\State Fit the predictor \(S_\psi\) on this teacher data using the declared loss family
\State For a new weighted cloud \((X,w)\), compute the student prediction \(h_\psi=S_\psi(X,w;\varepsilon,N,d_x)\)
\State Translate \(h_\psi\) into the initial state of the retained corrective Sinkhorn solver
\State Run \(K_{\mathrm{corr}}\) corrective Sinkhorn updates to obtain \(\widehat\Pi_{\psi,\varepsilon,K_{\mathrm{corr}}}\)
\State Form the barycentric equal-weight output cloud \(\widetilde X_{\psi,\mathrm{corr}}=B(\widehat\Pi_{\psi,\varepsilon,K_{\mathrm{corr}}})\)
\State Record row/column residuals, teacher-student discrepancy, and any downstream scalar checks
\end{algorithmic}
\end{algorithm}

\paragraph{What the student is allowed to change.}
Only the initialization of the retained teacher solver, or another explicitly
named latent teacher coordinate.

\paragraph{What the teacher still protects.}
The declared transport constraints, the EOT objective, and the corrected
barycentric output after refinement.

\section{Evidence contract for retained-teacher learning}
\label{sec:bf-learned-ot-training-objective}
\label{sec:bf-learned-ot-residuals}

Let \(h_\tau\) denote the teacher latent target and let
\(\mathcal D(N,d_x,\varepsilon,m,\eta)\) be the training distribution over
problem instances.  A simple supervised teacher-student fitting problem is
\begin{equation}
\label{eq:bf-learned-ot-training-objective}
  \psi^\star
  \in
  \arg\min_\psi
  \mathbb E_{(X,w,\varepsilon,m)\sim\mathcal D}
  \left[
    \mathcal L\{S_\psi(X,w;\varepsilon,N,d_x),h_\tau(X,w,\varepsilon)
    \}
  \right].
\end{equation}
This supervised loss is only the front end of the claim.  The chapter’s actual
deployment contract is about whether the corrected route still reproduces the
declared finite teacher closely enough on held-out and stress-test clouds.

The first deployment residual is therefore the teacher-student discrepancy after
correction:
\begin{equation}
\label{eq:bf-learned-ot-residual}
  R_{\psi,\tau}(X,w,\varepsilon)
  =
  B(\widehat\Pi_{\psi,\varepsilon,K_{\mathrm{corr}}})
  -
  B(\Pi_\tau(X,w,\varepsilon)).
\end{equation}
This residual does not by itself control filtering or HMC target error.  It only
measures how well the corrected student reproduces the chosen teacher.

A practical retained-teacher claim should therefore keep the following evidence
contract explicit:
\begin{enumerate}[label=(\arabic*)]
  \item the predicted latent teacher object is named explicitly,
  \item the corrective route satisfies the declared row/column residual contract,
  \item the corrected cloud discrepancy remains below the declared threshold on
  the tested envelope,
  \item the correction budget \(K_{\mathrm{corr}}\) or stopping rule is declared,
  \item any downstream filtering scalar or HMC-adjacent target is checked only
  after the corrected output is formed.
\end{enumerate}

This is also where the exact-vs-engineering distinction matters.  The exact
regularized teacher \(\Pi_\varepsilon^\star\), the finite numerical teacher
\(\Pi_{\varepsilon,K}\), and the learned warm start \(h_\psi\) are not the same
object.  The retained-teacher route is justified only when the engineering path
still returns to the declared corrected teacher closely enough for the stated
purpose.

\section{Permutation symmetry}
\label{sec:bf-learned-ot-symmetry}

Particle labels are arbitrary.  For a permutation matrix
\(P\in\mathbb R^{N\times N}\), the predictor should at least satisfy a declared
permutation symmetry.  In the output-cloud case the natural condition is
\begin{equation}
\label{eq:bf-learned-ot-equivariance}
  S_\psi(PX,Pw;\varepsilon,N,d_x)
  =
  P\,S_\psi(X,w;\varepsilon,N,d_x)
  \quad
  \text{or the corresponding latent-state analogue.}
\end{equation}
This is a necessary design condition, not a teacher-correctness theorem.

\section{Filtering-loop insertion point}
\label{sec:bf-learned-ot-loop}

Inside the filtering loop, the retained-teacher neural route sits exactly where
Chapter~\ref{ch:bf-dpf-entropic-ot-sinkhorn} placed the finite EOT/Sinkhorn
layer:
\begin{equation}
\label{eq:bf-learned-ot-loop}
  \{x_i,w_i\}_{i=1}^N
  \xrightarrow{\mathrm{student\ warm\ start + retained\ Sinkhorn}}
  \{\tilde x^{(j)},1/N\}_{j=1}^N.
\end{equation}
It accelerates the equal-weighting layer after the weighted cloud has already
been formed.  It does not redefine the earlier particle-weighting rule itself.

\section{Verification and non-claims}
\label{sec:bf-learned-ot-verification}
\label{sec:bf-learned-ot-boundary}

A practical implementation-facing reading of this chapter should still check:
\begin{itemize}
  \item teacher generation is declared and reproducible,
  \item the learned latent state is named and interpreted consistently,
  \item the retained corrective solve satisfies the declared coupling residuals,
  \item the corrected output cloud matches the declared finite teacher on the
  tested regime,
  \item downstream scalar checks are performed only on the corrected output when
  the route is embedded in filtering or HMC-adjacent code.
\end{itemize}

These checks are mathematical and numerical, not bureaucratic.  Their purpose is
simply to tell the reader whether the implementation still computes the declared
teacher-student object or has drifted to a different one.

This chapter does \emph{not} claim:
\begin{itemize}
  \item that the student replaces the teacher object,
  \item that small teacher-student residuals automatically imply posterior or HMC
  correctness,
  \item that broader neural OT families are unimportant, only that they answer a
  different implementation question.
\end{itemize}
Its role is narrower: fix one retained-teacher acceleration route in a form that
a reader can nearly implement directly.

\FloatBarrier
